% Abhakseite „Das kann ich“ – Zone nur im Gesamt (\mitzone)
%% TODO Ich-kann-Sätze: Platzhalter wie in den Aufgabendateien
\begin{abhakseite}
\ifdefined\mitzone
\abhakgruppe{Kennst du schon}
\abhak{1}{Binomialkoeffizient (n über k) als Anzahl der Auswahlen, Rechnertaste und kleine Werte im Kopf}
\abhak{2}{Pfadregeln ohne Zurücklegen (Bruchkette mit schrumpfendem Nenner) und das Gegenereignis}
\abhak{3}{Das Binomialmodell und seine Bedingungen (zur Abgrenzung: wann es ungeeignet ist)}
\abhak{4}{Brüche multiplizieren und kürzen}
\abhak{5}{Pfadregeln ohne Zurücklegen (Bruchkette mit schrumpfendem Nenner) und das Gegenereignis – Fehler finden}
\abhak{6}{Pfadregeln ohne Zurücklegen (Bruchkette mit schrumpfendem Nenner) und das Gegenereignis – selbst rechnen}
\fi
\abhakgruppe{Einheit 1 · Genau k Treffer ohne Zurücklegen: die Situation erkennen (feste kleine Gesamtheit, Ziehen ohne Zurücklegen}
\abhak{7}{Genau k ohne Zurücklegen}
\abhak{8}{Genau k ohne Zurücklegen – Fehler finden, Begründen, Anwendung}
\abhakgruppe{Einheit 2 · Kumulieren gegen eine Schranke: hypergeometrische Einzelwahrscheinlichkeiten aufsummieren und die größte (oder kleinste) Trefferzahl gegen eine Schranke bestimmen}
\abhak{9}{Kumulieren}
\abhak{10}{Kumulieren – Fehler finden, Begründen}
\end{abhakseite}
